%-----------------------------------------------------------------------
\appendix
\section{Material models}
\label{app:material}

Both media are evaluated from published measurement rather than from a fitted
dispersion curve; Section~\ref{sec:verify-index} records why the fit is the wrong
instrument here.

\subsection{Water}

Equation~\eqref{eq:sellmeier} with the four-term Sellmeier coefficients of
Daimon and Masumura \cite{daimon2007}, measured on HPLC-distilled water at
\SI{20.0}{\celsius}. The coefficients as evaluated by the solver are:

\begin{table}[htbp]
\centering
\caption{Water Sellmeier coefficients, $\lambda$ in micrometres. These are the
values in \texttt{src/app/showcase\_main.cpp} and in the Python reference, and
they agree term for term.}
\label{tab:sellmeier}
\small
\begin{tabular}{@{}rr@{}}
\toprule
$j$ & $B_j$ \hspace{3em} $C_j$ \\
\midrule
1 & \num{5.684027565e-1} \hspace{3em} \num{5.101829712e-3} \\
2 & \num{1.726177391e-1} \hspace{3em} \num{1.821153936e-2} \\
3 & \num{2.086189578e-2} \hspace{3em} \num{2.620722293e-2} \\
4 & \num{1.130748688e-1} \hspace{3em} \num{1.069792721e1} \\
\bottomrule
\end{tabular}
\end{table}

\subsection{Glass}

Equation~\eqref{eq:rubin} with the soda-lime silica coefficients of Rubin
\cite{rubin1985}: $1.5130 - 0.003169\,\lambda^2 + 0.003962/\lambda^2$ with
$\lambda$ in micrometres.

\subsection{Evaluated indices}

\begin{table}[htbp]
\centering
\caption{Indices at the three render channels. Red has the lowest index in both
media, so red refracts least and lands farthest from the axis; the chromatic
ordering in every figure follows from this table, and from nothing else.}
\label{tab:indices}
\small
\begin{tabular}{@{}lrrr@{}}
\toprule
Channel & $\lambda$ & Water $n$ & Glass $n$ \\
\midrule
R & \SI{650}{\nano\meter} & \num{1.33167} & \num{1.52104} \\
G & \SI{550}{\nano\meter} & \num{1.33468} & \num{1.52514} \\
B & \SI{450}{\nano\meter} & \num{1.33961} & \num{1.53192} \\
\bottomrule
\end{tabular}
\end{table}

\subsection{The withdrawn fit}

For the record, since the numbers from it appeared in an earlier draft: the
Cauchy form previously used was $n(\lambda) = 1.3208 + 5200/\lambda^2$ with
$\lambda$ in nanometres for water, and $1.5095 + 4200/\lambda^2$ for glass.
Evaluated at \SI{450}{\nano\meter} and \SI{650}{\nano\meter}, the water form
carries roughly \num{1.97} times the dispersion of \eqref{eq:sellmeier}, and the
glass form is about \num{0.0017} low at one channel. These forms must not be
revived; the measured path is the one in Table~\ref{tab:indices}.

%-----------------------------------------------------------------------
\section{The fold normal form}
\label{app:normalform}

Proposition~\ref{prop:foldnormal} is stated for a one-dimensional pupil. Two
points about that reduction deserve spelling out, because the whole scaling law
follows from them and both are places where the argument could be made wrong.

\paragraph{Why one dimension is enough.} The map \eqref{eq:raymap} is
two-dimensional, but at a fold the Jacobian has rank one, so there is a single
direction in which the map is degenerate and a transverse direction in which it
is a local diffeomorphism. In the normal-form coordinates $(u,v)$ of
\eqref{eq:normalform}, $u$ is the non-degenerate direction: it takes part in no
degeneracy, and the diffraction integral factors into an integral along $u$
times the one-dimensional integral \eqref{eq:phase} along $v$. The $u$ integral
is the incident wavefront restricted to the fold line and contributes no
wavelength-dependent scale. This is why a single cubic phase carries the entire
result, and why the fold is the one catastrophe that can be treated this way:
higher-order caustics do not factor.

\paragraph{Why the pupil truncation does not matter.} The integral
\eqref{eq:phase} is stated over $\rho \in [-1,1]$, a finite pupil, while the
Airy function is an integral over the whole line. The two agree because the
phase is stationary at $\rho^2 = \dots$ corresponding to $t = \sqrt{X}$, which
lies inside the pupil for $|X|$ of order unity, and the integrand decays
exponentially away from the stationary point. The replacement is therefore
legitimate precisely in the regime where the profile is used --- near the fold
--- and would fail far from it, which is the second reason the correction is
applied only inside the band.

\paragraph{The scaling, assembled.} Writing the canonical coordinate from the
proof of Proposition~\ref{prop:foldnormal} as
\begin{equation}
X = -x\,k^{2/3}2^{1/3}c^{-1/3},
\end{equation}
the transverse distance at which $|X| = 1$ is
\begin{equation}
d = k^{-2/3}2^{-1/3}c^{1/3}
  = \frac{\lambda^{2/3}c^{1/3}}{2\pi^{2/3}},
\end{equation}
which contains $\lambda$ only as $\lambda^{2/3}$. Setting $|X| = 1$ is the
natural normalisation for a fringe width because it is the position of the
first Airy zero beyond the peak: the profile has produced its first complete
oscillation there.

%-----------------------------------------------------------------------
\section{Estimators, formally}
\label{app:estimators}

Section~\ref{sec:estimators} lists the statistics informally. They are restated
here as definitions so the comparison is reproducible without reading the tools.

\begin{definition}[Edge radius]
\label{def:edge}
For a single-channel image $I$ on a receiver with known centre $\mathbf{0}$,
let $M = \max I$ and, for $r > 0$, let $\bar I(r)$ be the mean of $I$ over the
annulus $\{\,r \le \|\mathbf{u}\| < r + \SI{1}{\milli\meter}\,\}$. The edge
radius is
\begin{equation}
r_e \;=\; \sup\{\, r : \bar I(r) > 0.02\,M \,\}.
\end{equation}
Dividing by $M$ makes the statistic invariant to any uniform scale, and taking a
mean over an annulus rather than a maximum makes it unaffected by a ceiling
applied to the core.
\end{definition}

\begin{definition}[Tiled correlation]
\label{def:corr}
For images $A$, $B$ on a common grid and tile edge $s$ in pixels, partition the
grid into $\lceil W/s \rceil \times \lceil H/s \rceil$ tiles, replace each tile by
its mean, and compute the Pearson correlation of the resulting two matrices. The
statistic is reported by tile edge $s$; naming it by tile \emph{count} instead
reverses the apparent trend, which is the error of
Section~\ref{sec:verify-metric}.
\end{definition}

\begin{definition}[Positive-part differential]
\label{def:pospart}
Given renders $P$ (with the refracting body present) and $Q$ (without), the
caustic estimate is $\max(P - Q, 0)$ elementwise. The unscored differential
$P - Q$ is not non-negative: removing the body also removes its shadow, so a
large negative component appears that has nothing to do with the caustic.
\end{definition}

\begin{definition}[Temporal flicker]
\label{def:flicker}
For a rendered sequence $\{I_t\}$ and a region $R$, the flicker is
\begin{equation}
F \;=\; \frac{1}{|R|}\sum_{\mathbf{u} \in R}
\frac{1}{T-1}\sum_{t=1}^{T-1}\big| I_{t+1}(\mathbf{u}) - I_t(\mathbf{u}) \big| ,
\end{equation}
in units of the 8-bit quantisation step.
\end{definition}

%-----------------------------------------------------------------------
\section{Notation and reproducibility manifest}
\label{app:notation}
\label{app:manifest}

\begin{table}[htbp]
\centering
\caption{Symbols used throughout.}
\label{tab:notation}
\footnotesize
\begin{tabular}{@{}ll@{}}
\toprule
Symbol & Meaning \\
\midrule
$\mathbf{x}, h, \nrm$ & surface coordinate, height field \eqref{eq:surface}, normal \eqref{eq:normal} \\
$\dir, \trn, \eta$ & source direction, refracted direction \eqref{eq:snell}, index ratio \\
$c_i, c_t, T$ & cosines of incidence/transmission, Fresnel power \eqref{eq:fresnel} \\
$\Hm, \Jm, \detJ$ & ray map \eqref{eq:raymap}, Jacobian, and determinant \\
$\fold, \recv, I$ & fold locus $\{\detJ = 0\}$, receiver footprint, irradiance \eqref{eq:irradiance} \\
$B, s, w_c$ & detected fold band, signed distance, fringe width \eqref{eq:perchannel} \\
$\Ai, X, b(s)$ & Airy function, canonical coordinate \eqref{eq:canonical}, profile \eqref{eq:profile} \\
$\kappa, k^{*}$ & display ceiling, band flux ratio \eqref{eq:kstar} \\
\bottomrule
\end{tabular}
\end{table}

\begin{table}[htbp]
\centering
\caption{Bake parameters used for all results in this paper.}
\label{tab:manifest}
\footnotesize
\begin{tabular}{@{}ll@{}}
\toprule
Parameter & Value \\
\midrule
Bake grid (results / study) & $1024^2$ / $256^2$--$1024^2$ \\
Spatial extent \& receiver  & \SI{\pm 160}{\milli\meter}, \SI{320}{\milli\meter} square patch \\
Surface amplitudes $A_i$    & \num{0.65}, \num{0.40}, \num{0.25}, \num{0.15} \si{\milli\meter} \\
Surface wavelengths $\Lambda_i$ & 42, 27, 17, 11 \si{\milli\meter} \\
Source & distant, \SI{0.53}{\degree}, direction \eqref{eq:snell} \\
Detector threshold & 2nd percentile of $|\detJ|$ \\
Reference width $w_{\text{ref}}$ & \SI{123}{\micro\meter} (physical), 2048-LUT on $[0,8]$ \\
Profile floor \& ceiling & floor \num{0.15}, ceiling $1.35\,\kappa$ \\
\bottomrule
\end{tabular}
\end{table}

The reference width is the physical fringe scale for this geometry, obtained
independently by \texttt{tools/fringe\_scale.py}. It is deliberately not
inflated for visibility: an earlier version of this work used \SI{1.5}{}
\si{\milli\meter}, which is \num{12.2} times the physical value, and the
difference between the two is the difference between showing a measured
dispersion and drawing decorative fringes.
